\scanpage{76}
\mspage{37}
9) \uline{La structure groupoïdale des \inserted{\uline{espaces modulaires de}} Teichmüller \unsure{modules|variétés}} (la \unsure{MDT|MDS}-structure) : cas transcendant, morphismes de générisation.

Nous travaillons maintenant avec les $M^{\mathrm{an}}_{g,\nu}$, $\widehat M^{\mathrm{an}}_{g,\nu}$ etc -- et pour ne pas traîner les exposants ``an'' (pour ``analytique''), nous les faisons tomber.

Considérons pour tout MD-graphe $\underline G$, la multiplicité modulaire (analytique) $M_{\underline G}$, et son groupoïde fondamental, que je note $\Pi M_{\underline G}$. C'est une \inserted{``grosse''} catégorie, par la suite nous allons voir comment la remplacer \inserted{\unsure{canoniqt}} par une \deleted{catégorie} \inserted{groupoïde} dont l'ens des objets soit fini. Dans nos constructions, il est en effet licite de remplacer un groupoïde tel que $\Pi M_{\underline G}$ par un groupoïde équivalent, pourvu que \deleted{l'équivalence} \inserted{dans celui-ci} \deleted{\unsure{un groupe}} \inserted{le groupe}
\[
(94)\qquad \Gamma_{\underline G}=\operatorname{Aut}\underline G
\]
opère, et que l'équivalence commute à l'action de ce groupe.

Notons déjà que $\Pi M_{\underline G}$ est un groupoïde
\margin{75}
